% Alternation of process timeslices and scheduler runs.
% Usage: % Alternation of process timeslices and scheduler runs.
% Usage: % Alternation of process timeslices and scheduler runs.
% Usage: % Alternation of process timeslices and scheduler runs.
% Usage: \input{figures/l02-timeslice}   (width ~117 mm incl. the CPU label)
\begin{tikzpicture}[x=1mm, y=1mm, font=\footnotesize\sffamily,
  dim/.style={{Bar[width=2mm]}-{Bar[width=2mm]}, thin},
  pr/.style={draw, fill=ln-accent!22},
  sc/.style={draw, fill=ln-box}]
  \node[anchor=east] at (-1,3) {CPU};
  \draw[pr] (0,0)  rectangle (30,6)  node[midway] {P1};
  \draw[sc] (30,0) rectangle (38,6)  node[midway, font=\scriptsize\sffamily] {\iflangja{S}{S}};
  \draw[pr] (38,0) rectangle (68,6)  node[midway] {P2};
  \draw[sc] (68,0) rectangle (76,6)  node[midway, font=\scriptsize\sffamily] {\iflangja{S}{S}};
  \draw[pr] (76,0) rectangle (106,6) node[midway] {P3};
  \draw[sc] (106,0) rectangle (109,6);
  \foreach \a/\b in {0/30, 38/68, 76/106} {
    \draw[dim] (\a,-3) -- (\b,-3);
    \node[below] at ({(\a+\b)/2},-3.5) {$T_p$};
  }
  \foreach \a/\b in {30/38, 68/76} {
    \draw[dim] (\a,9) -- (\b,9);
    \node[above] at ({(\a+\b)/2},9.5) {$t_{\mathrm{sched}}$};
  }
\end{tikzpicture}
   (width ~117 mm incl. the CPU label)
\begin{tikzpicture}[x=1mm, y=1mm, font=\footnotesize\sffamily,
  dim/.style={{Bar[width=2mm]}-{Bar[width=2mm]}, thin},
  pr/.style={draw, fill=ln-accent!22},
  sc/.style={draw, fill=ln-box}]
  \node[anchor=east] at (-1,3) {CPU};
  \draw[pr] (0,0)  rectangle (30,6)  node[midway] {P1};
  \draw[sc] (30,0) rectangle (38,6)  node[midway, font=\scriptsize\sffamily] {\iflangja{S}{S}};
  \draw[pr] (38,0) rectangle (68,6)  node[midway] {P2};
  \draw[sc] (68,0) rectangle (76,6)  node[midway, font=\scriptsize\sffamily] {\iflangja{S}{S}};
  \draw[pr] (76,0) rectangle (106,6) node[midway] {P3};
  \draw[sc] (106,0) rectangle (109,6);
  \foreach \a/\b in {0/30, 38/68, 76/106} {
    \draw[dim] (\a,-3) -- (\b,-3);
    \node[below] at ({(\a+\b)/2},-3.5) {$T_p$};
  }
  \foreach \a/\b in {30/38, 68/76} {
    \draw[dim] (\a,9) -- (\b,9);
    \node[above] at ({(\a+\b)/2},9.5) {$t_{\mathrm{sched}}$};
  }
\end{tikzpicture}
   (width ~117 mm incl. the CPU label)
\begin{tikzpicture}[x=1mm, y=1mm, font=\footnotesize\sffamily,
  dim/.style={{Bar[width=2mm]}-{Bar[width=2mm]}, thin},
  pr/.style={draw, fill=ln-accent!22},
  sc/.style={draw, fill=ln-box}]
  \node[anchor=east] at (-1,3) {CPU};
  \draw[pr] (0,0)  rectangle (30,6)  node[midway] {P1};
  \draw[sc] (30,0) rectangle (38,6)  node[midway, font=\scriptsize\sffamily] {\iflangja{S}{S}};
  \draw[pr] (38,0) rectangle (68,6)  node[midway] {P2};
  \draw[sc] (68,0) rectangle (76,6)  node[midway, font=\scriptsize\sffamily] {\iflangja{S}{S}};
  \draw[pr] (76,0) rectangle (106,6) node[midway] {P3};
  \draw[sc] (106,0) rectangle (109,6);
  \foreach \a/\b in {0/30, 38/68, 76/106} {
    \draw[dim] (\a,-3) -- (\b,-3);
    \node[below] at ({(\a+\b)/2},-3.5) {$T_p$};
  }
  \foreach \a/\b in {30/38, 68/76} {
    \draw[dim] (\a,9) -- (\b,9);
    \node[above] at ({(\a+\b)/2},9.5) {$t_{\mathrm{sched}}$};
  }
\end{tikzpicture}
   (width ~117 mm incl. the CPU label)
\begin{tikzpicture}[x=1mm, y=1mm, font=\footnotesize\sffamily,
  dim/.style={{Bar[width=2mm]}-{Bar[width=2mm]}, thin},
  pr/.style={draw, fill=ln-accent!22},
  sc/.style={draw, fill=ln-box}]
  \node[anchor=east] at (-1,3) {CPU};
  \draw[pr] (0,0)  rectangle (30,6)  node[midway] {P1};
  \draw[sc] (30,0) rectangle (38,6)  node[midway, font=\scriptsize\sffamily] {\iflangja{S}{S}};
  \draw[pr] (38,0) rectangle (68,6)  node[midway] {P2};
  \draw[sc] (68,0) rectangle (76,6)  node[midway, font=\scriptsize\sffamily] {\iflangja{S}{S}};
  \draw[pr] (76,0) rectangle (106,6) node[midway] {P3};
  \draw[sc] (106,0) rectangle (109,6);
  \foreach \a/\b in {0/30, 38/68, 76/106} {
    \draw[dim] (\a,-3) -- (\b,-3);
    \node[below] at ({(\a+\b)/2},-3.5) {$T_p$};
  }
  \foreach \a/\b in {30/38, 68/76} {
    \draw[dim] (\a,9) -- (\b,9);
    \node[above] at ({(\a+\b)/2},9.5) {$t_{\mathrm{sched}}$};
  }
\end{tikzpicture}
